\documentclass[a4paper,11pt]{article}
\usepackage{exam-style-341}
\examinehead{341 农业综合知识三 · 数据库分章练习}{第05章\quad 事务、并发与恢复\quad 参考答案}
\begin{document}

\unitanswerscover{341 农业综合知识三 · 数据库分章练习}{第05章\quad 事务、并发与恢复 参考答案}{答案与解析按练习册题号顺序编排；判断题先给出判断，再给出改正后的正确表述；综合题给出 SQL 事务框架、并发问题判定与日志恢复的完整推演过程。}

\answerpart{一、单项选择题（10 题，每题 2 分，共 20 分）}

\q \textbf{D}。\begin{analysis}持久性（Durability）指事务一旦提交，它对数据库中数据的改变就是永久性的，后续的其他操作或系统故障不应影响其结果。原子性强调“全做或全不做”；一致性强调状态转换的正确性；隔离性强调并发事务互不干扰。\end{analysis}

\q \textbf{C}。\begin{analysis}读脏数据是指事务读到了另一个事务修改但尚未提交的数据，而该事务随后被回滚，所读数据与数据库中的实际数据不一致。本题中 $T_2$ 读到的 900 是 $T_1$ 未提交且被回滚的值，属于读脏数据。\end{analysis}

\q \textbf{C}。\begin{analysis}排他锁（X 锁）与任何其他锁都不相容：一个事务对数据加了 X 锁后，其他事务不能对该数据加任何锁（既不能加 X 锁，也不能加 S 锁），只能等待锁释放。共享锁（S 锁）之间相容。\end{analysis}

\q \textbf{A}。\begin{analysis}一级封锁协议规定事务在修改数据前必须先加 X 锁并保持到事务结束，可以防止丢失修改。一级协议不能防止读脏数据和不可重复读；二级协议进一步防止读脏数据；三级协议进一步防止不可重复读。\end{analysis}

\q \textbf{A}。\begin{analysis}两段锁协议把事务分为两个阶段：扩展阶段只能加锁不能解锁，收缩阶段只能解锁不能加锁。遵守两段锁协议的并发调度一定是可串行化的。\end{analysis}

\q \textbf{A}。\begin{analysis}死锁的预防方法包括一次封锁法和顺序封锁法：一次封锁法要求事务一次性地封锁所需全部数据；顺序封锁法要求所有事务按同一顺序申请封锁。超时法和事务等待图法是死锁的诊断方法，检查点法用于恢复。\end{analysis}

\q \textbf{A}。\begin{analysis}事务故障是指某个事务未能正常结束。恢复子系统反向扫描日志文件，对该事务的更新执行 UNDO，将其修改恢复到修改前的值。REDO 用于系统故障后处理已提交事务；重装副本用于介质故障。\end{analysis}

\q \textbf{B}。\begin{analysis}系统故障恢复采用“先 UNDO 后 REDO”的策略：对故障发生时尚未提交的事务执行 UNDO，对已提交但修改可能仍留在缓冲区而未写回磁盘的事务执行 REDO。因此已提交事务应执行 REDO。\end{analysis}

\q \textbf{A}。\begin{analysis}检查点的作用是：在检查点时刻强制把日志记录和修改过的数据写回磁盘，并记录此时正在执行的事务。恢复时只需从检查点开始扫描日志，对检查点之后提交的事务执行 REDO、对未提交的事务执行 UNDO，从而减少需要扫描的日志范围和重做的事务，缩短恢复时间。\end{analysis}

\q \textbf{A}。\begin{analysis}GRANT 和 REVOKE 是自主存取控制中的授权与收回权限语句，属于数据库安全性控制。完整性控制针对数据的正确性和相容性；并发控制针对事务之间的干扰；恢复控制针对故障后的数据恢复。\end{analysis}

\answerpart{二、多项选择题（5 题，每题 4 分，共 20 分）}

\q \textbf{ABCD}。\begin{analysis}ACID 即原子性（Atomicity）、一致性（Consistency）、隔离性（Isolation）和持久性（Durability），四项都是事务的基本性质。\end{analysis}

\q \textbf{ABC}。\begin{analysis}并发操作可能带来三类数据不一致性问题：丢失修改、不可重复读和读脏数据。D 数据冗余属于数据库设计问题，不是并发操作带来的不一致性。\end{analysis}

\q \textbf{ACD}。\begin{analysis}共享锁之间相容，排他锁与共享锁不相容，排他锁与排他锁也不相容。B 错误：共享锁与排他锁不相容。\end{analysis}

\q \textbf{ABC}。\begin{analysis}数据库系统通常将故障分为事务故障、系统故障和介质故障三类。D 计算机病毒可能造成数据破坏，但不是数据库恢复理论中的标准故障分类。\end{analysis}

\q \textbf{ABCD}。\begin{analysis}日志记录事务对数据库的每一次更新操作，通常包含事务标识、数据项标识、修改前的值、修改后的值，以及事务开始、提交、回滚等标记；利用日志可以执行 UNDO 和 REDO；为保证可恢复性，必须先写日志、后写数据库（WAL 原则）。四项均正确。\end{analysis}

\answerpart{三、判断题（5 题，每题 2 分，共 10 分）}

\q √。\begin{analysis}原子性要求事务中的操作要么全部执行，要么全部不执行；若事务中途失败，已做的修改必须全部撤销，使数据库回到事务开始前的状态。\end{analysis}

\q √。\begin{analysis}两段锁协议要求所有事务分扩展阶段和收缩阶段加锁、解锁：扩展阶段只能加锁不能解锁，收缩阶段只能解锁不能加锁。可以证明，遵守两段锁协议的并发调度一定是冲突可串行化的。注意两段锁协议是充分条件而非必要条件，且可能产生死锁。\end{analysis}

\q ×。改正：死锁属于事务故障，不是介质故障。\begin{analysis}当两个或多个事务互相等待对方持有的锁时产生死锁，DBMS 通常选中其中一个事务撤销，使其回滚，这属于事务故障。介质故障是磁盘损坏等造成的数据丢失。\end{analysis}

\q ×。改正：系统故障恢复时，对故障发生前已提交的事务应执行 REDO，对尚未提交的事务执行 UNDO。\begin{analysis}已提交事务的修改可能仍在缓冲区中而未写回磁盘，系统故障后内存信息丢失，因此必须依据日志执行 REDO，重新写入其结果；未提交事务的修改可能已部分写回磁盘，必须执行 UNDO 撤销。\end{analysis}

\q ×。改正：读脏数据是指事务读到了另一个事务修改但尚未提交的数据，而该事务随后被回滚。\begin{analysis}若读到的是已提交的数据，其值已得到确认，不属于“脏”数据。读脏数据的危害在于所读数据可能因回滚而消失或与数据库实际值不一致。\end{analysis}

\answerpart{四、填空题（5 题，每题 2 分，共 10 分）}

\q 第一空：\textbf{一致性}；第二空：\textbf{持久性}。\begin{analysis}ACID 即原子性（Atomicity）、一致性（Consistency）、隔离性（Isolation）、持久性（Durability）。\end{analysis}

\q 第一空：\textbf{不可重复读}；第二空：\textbf{读脏数据}。\begin{analysis}并发操作带来的三类不一致性问题是丢失修改、不可重复读和读脏数据。\end{analysis}

\q 第一空：\textbf{排他（X）}；第二空：\textbf{共享（S）}。\begin{analysis}排他锁又称写锁（X 锁），共享锁又称读锁（S 锁）。加 X 锁后其他事务不能加任何锁；加 S 锁后其他事务只能再加 S 锁。\end{analysis}

\q 第一空：\textbf{丢失修改}；第二空：\textbf{读脏数据}。\begin{analysis}一级封锁协议要求修改前加 X 锁并保持到事务结束，防止丢失修改；二级封锁协议在此基础上要求读前加 S 锁、读完即释放，进一步防止读脏数据；三级封锁协议要求读锁保持到事务结束，进一步防止不可重复读。\end{analysis}

\q 第一空：\textbf{系统故障}；第二空：\textbf{介质故障}。\begin{analysis}数据库系统的故障分为事务故障、系统故障和介质故障三类，恢复方法分别以 UNDO、UNDO+REDO、重装副本+REDO 为主。\end{analysis}

\answerpart{五、简答题（2 题，每题 8 分，共 16 分）}

\q 参考答案要点：

\textbf{（1）事务的 ACID 性质（4 分）}
（1）原子性：事务中的操作要么全部执行，要么全部不执行；若中途失败，已做的修改全部撤销。
（2）一致性：事务执行的结果必须使数据库从一个一致性状态变到另一个一致性状态。
（3）隔离性：一个事务的执行不能被其他事务干扰，并发事务之间不能互相看到未提交的中间结果。
（4）持久性：事务一旦提交，它对数据库中数据的改变就是永久性的。

\textbf{（2）三类数据不一致性问题（4 分）}
（1）丢失修改：两个事务读入同一数据并各自修改，后提交的事务覆盖了先提交事务的修改，使先前的修改丢失。
（2）不可重复读：事务先后两次读取同一数据，由于其他事务在两次读之间进行了修改、插入或删除，两次结果不一致。
（3）读脏数据：事务读到另一个事务修改但尚未提交的数据，而该事务随后被回滚，所读数据是错误的。
\begin{analysis}评分要点：ACID 四项各 1 分；三类并发问题各 1 分，且能说明产生原因或现象再得 1 分。常见错误：把“不可重复读”和“读脏数据”混淆；漏写“事务未提交”这一关键限定。\end{analysis}

\q 参考答案要点：

\textbf{（1）一级封锁协议（2.5 分）}
事务在修改数据 $R$ 之前必须先对其加 X 锁，直到事务结束才释放。一级封锁协议可以防止丢失修改，但不能防止读脏数据和不可重复读。

\textbf{（2）二级封锁协议（2.5 分）}
在一级封锁协议的基础上，事务在读取数据 $R$ 之前必须先对其加 S 锁，读完即可释放 S 锁。二级封锁协议可以防止丢失修改和读脏数据，但不能防止不可重复读。

\textbf{（3）三级封锁协议（2.5 分）}
在一级封锁协议的基础上，事务在读取数据 $R$ 之前必须先对其加 S 锁，直到事务结束才释放。三级封锁协议可以防止丢失修改、读脏数据和不可重复读。

\textbf{（4）归纳（0.5 分）}
三级封锁协议逐级加强：级别越高，能避免的并发问题越多，但封锁时间越长，并发度越低。
\begin{analysis}评分要点：三级协议各自的加锁规则（各 1.5 分）、各自能解决的问题（各 1 分）、逐级加强的归纳（0.5 分）。常见错误：把二级、三级协议中 S 锁的释放时机写反；误认为一级协议能防止不可重复读。\end{analysis}

\answerpart{六、综合题（2 题，每题 12 分，共 24 分）}

\q 参考答案：

\textbf{（1）SQL 事务框架（3 分）}
\begin{lstlisting}[language=SQL]
START TRANSACTION;

UPDATE Account
SET Balance = Balance - 100
WHERE Ano = 'A';

UPDATE Account
SET Balance = Balance + 100
WHERE Ano = 'B';

COMMIT;
\end{lstlisting}
说明：两条 UPDATE 必须放在同一个事务中。若第二条更新失败或业务校验不通过，应执行 \code{ROLLBACK}，撤销第一条 UPDATE 已做的扣款，从而保证“要么都做、要么都不做”的原子性。MySQL 中 \code{START TRANSACTION} 显式开启事务，\code{COMMIT} 提交，\code{ROLLBACK} 回滚；也可用 \code{SET autocommit=0} 关闭自动提交。

\textbf{（2）三个并发场景的判断（6 分）}

场景一：\textbf{丢失修改}。$T_1$ 和 $T_2$ 都读到 $A=1000$ 并各自修改后写回，$T_2$ 写回的 800 覆盖了 $T_1$ 写回的 900，$T_1$ 的修改被丢失。

场景二：\textbf{不可重复读}。$T_1$ 第一次读 $A$ 得 1000，$T_2$ 把 $A$ 改为 900 并提交，$T_1$ 再次读 $A$ 得 900，同一事务内两次读取结果不同。

场景三：\textbf{读脏数据}。$T_1$ 把 $A$ 改为 900 但尚未提交，$T_2$ 读到 900；随后 $T_1$ 回滚，$A$ 恢复为 1000，$T_2$ 读到的 900 是错误数据。

\textbf{（3）两段锁协议（3 分）}
两段锁协议要求所有事务分两个阶段对数据项加锁和解锁：
（1）扩展阶段：事务可以申请锁，但不能释放任何锁；
（2）收缩阶段：事务可以释放锁，但不能申请新的锁。
两段锁协议的作用是保证并发调度的可串行化：遵守两段锁协议的并发调度一定是冲突可串行化的，从而保证并发执行结果的正确性。但两段锁协议不能避免死锁，且它不是可串行化的必要条件。
\begin{analysis}评分要点：（1）事务控制语句完整、两条 UPDATE 位于同一事务（3 分）；（2）三个场景判断正确并说明理由（每个 2 分，共 6 分）；（3）两段锁协议的阶段划分与可串行化结论（3 分）。常见错误：把场景一判为“不可重复读”（其核心是后写覆盖前写导致丢失修改）；把场景三判为“不可重复读”（其关键是被读数据未提交且后来回滚）；把两段锁协议的两个阶段写反。\end{analysis}

\q 参考答案：

\textbf{（1）各事务在故障发生时的状态（3 分）}
\begin{itemize}[leftmargin=1.6em, itemsep=1pt]
  \item T1：已提交，且提交发生在检查点之前，其修改已由检查点保证写回磁盘；
  \item T2：已提交，但提交发生在检查点之后；
  \item T3：未提交（故障时仍在执行），属于未完成事务；
  \item T4：已提交，且提交发生在检查点之后。
\end{itemize}

\textbf{（2）UNDO 与 REDO 事务集合（3 分）}
\begin{itemize}[leftmargin=1.6em, itemsep=1pt]
  \item UNDO 集合：$\{$T3$\}$。T3 在故障时尚未提交，必须撤销其对 C 的修改。
  \item REDO 集合：$\{$T2, T4$\}$。T2、T4 在检查点之后提交，其修改可能仍在缓冲区而未写回磁盘，必须重做。
  \item T1 在检查点之前已提交，检查点已保证其修改写回磁盘，不需要 REDO。
\end{itemize}

\textbf{（3）恢复后各数据项的值（3 分）}
\begin{itemize}[leftmargin=1.6em, itemsep=1pt]
  \item $A=200$：T1 已提交且检查点前已写回，保持 200；
  \item $B=80$：对 T2 执行 REDO，$B$ 由 50 改为 80；
  \item $C=30$：对 T3 执行 UNDO，$C$ 由 60 恢复为修改前的 30；
  \item $D=20$：对 T4 执行 REDO，$D$ 由 10 改为 20。
\end{itemize}

\textbf{（4）检查点的作用（3 分）}
检查点时刻，DBMS 强制把日志缓冲区中的日志记录和数据库缓冲区中修改过的数据写回磁盘，并记录此时正在执行的事务。恢复时只需从检查点开始扫描日志：对检查点之后提交的事务执行 REDO，对故障时尚未提交的事务执行 UNDO；检查点之前已提交的事务不必重做。这样既保证了恢复的正确性，又减少了需要扫描的日志量和重做的事务数，缩短了恢复时间。
\begin{analysis}评分要点：四个事务状态判断（3 分）；UNDO、REDO 集合正确（3 分，漏掉或多出事务各扣分）；四个数据项的恢复值及依据（3 分）；检查点作用（3 分）。常见错误：（1）把检查点之前提交的 T1 也列入 REDO；（2）把已提交的 T2、T4 列入 UNDO；（3）认为故障后 C 仍为 60（忘记 T3 未提交需撤销）；（4）把检查点的作用说成“提高查询速度”或“防止死锁”。\end{analysis}

\end{document}
